\chapter{Diabetic Retinopathy}

\section*{Definition and Overview}
Diabetic retinopathy is damage to the retina, the light-sensitive tissue at the back of the eye, caused by chronically high blood glucose levels in people with diabetes. High glucose damages the small blood vessels of the retina, causing them to leak fluid and blood, become blocked, or grow abnormally. Retinopathy is a leading cause of preventable blindness in working-age adults, but it develops gradually, causes no symptoms in its early stages, and can be detected and treated before sight is threatened. Regular eye screening is essential for everyone with diabetes.

\section*{Epidemiology}
Diabetic retinopathy affects around one in three people with diabetes, and its prevalence rises with the duration of diabetes: after 20 years, most people with type 1 diabetes and a large proportion with type 2 diabetes have some retinopathy. It is the most common cause of blindness in Australians of working age. Risk is higher with poor glucose control, high blood pressure, long duration of diabetes, pregnancy, and smoking. Australia has a national diabetic retinopathy screening program in many states, which has helped detect disease earlier.

\section*{Aetiology and Pathophysiology}
\begin{itemize}
    \item \textbf{High blood glucose:} chronic hyperglycaemia damages the lining of retinal blood vessels
    \item \textbf{Non-proliferative retinopathy:} early changes with microaneurysms, small haemorrhages, and leaking vessels
    \item \textbf{Macular oedema:} fluid accumulates in the macula, the central part of the retina responsible for sharp vision
    \item \textbf{Proliferative retinopathy:} severe vessel blockage triggers growth of new, fragile vessels that bleed easily and can detach the retina
    \item \textbf{Risk factors:} duration of diabetes, poor glucose control, hypertension, pregnancy, and smoking
    \item \textbf{Sight loss:} occurs mainly through macular oedema and bleeding from new vessels
\end{itemize}

\section*{Clinical Features}
\begin{itemize}
    \item Early retinopathy causes no symptoms, which is why screening is vital
    \item Blurred or fluctuating vision as the macula becomes involved
    \item Floaters, spots, or cobweb-like shapes in the vision from bleeding into the vitreous
    \item Dark or empty areas in the vision
    \item Sudden, severe vision loss if there is major bleeding or retinal detachment
    \item Difficulty reading, driving, or recognising faces in advanced disease
\end{itemize}

\section*{Differential Diagnosis}
\begin{itemize}
    \item Other causes of retinal bleeding, such as retinal vein occlusion and hypertension-related retinopathy
    \item Age-related macular degeneration, which also affects the central retina
    \item Cataracts and glaucoma, other common causes of vision loss in people with diabetes
    \item Vitreous detachment and other causes of floaters
    \item Refractive changes, since blood glucose changes can temporarily alter vision
\end{itemize}

\section*{When to Seek Medical Care}
Everyone with diabetes should have an eye examination through dilated pupils or retinal photography at least every one to two years, or more often if retinopathy is found. Seek urgent eye care for any sudden change in vision, new floaters, or loss of sight.

\textbf{Call 000 (or local emergency services) for:}
\begin{itemize}
    \item Sudden loss of vision or a sudden shower of floaters and flashing lights
    \item A dark curtain or shadow over part of the vision
    \item Severe eye pain with redness
    \item Sudden weakness or collapse, which may indicate stroke
\end{itemize}

\section*{Diagnosis and Investigations}
\begin{itemize}
    \item \textbf{Dilated eye examination:} an ophthalmologist or optometrist examines the retina after dilating the pupils
    \item \textbf{Retinal photography:} digital images of the retina are used in screening programs and can be reviewed remotely
    \item \textbf{Optical coherence tomography (OCT):} detailed images of the retina that detect macular oedema
    \item \textbf{Fluorescein angiography:} a dye test that shows leaking and abnormal vessels
    \item \textbf{Glucose and blood pressure assessment:} levels guide treatment
\end{itemize}

\section*{Management}
\begin{itemize}
    \item \textbf{Glucose control:} the cornerstone of prevention and treatment; tight control slows progression
    \item \textbf{Blood pressure control:} lowers the risk of progression and macular oedema
    \item \textbf{Monitoring:} regular eye reviews with frequency based on the severity of retinopathy
    \item \textbf{Laser therapy:} photocoagulation seals leaking vessels and treats abnormal new vessels
    \item \textbf{Injections into the eye:} anti-VEGF medicines such as ranibizumab and aflibercept treat macular oedema and proliferative disease
    \item \textbf{Vitrectomy:} surgery to remove blood and scar tissue in advanced disease
    \item \textbf{Treatment of associated problems:} management of high cholesterol and smoking cessation
\end{itemize}

\section*{Complications}
\begin{itemize}
    \item Macular oedema causing central vision loss
    \item Vitreous haemorrhage from fragile new vessels
    \item Tractional retinal detachment
    \item Neovascular glaucoma, a serious complication of proliferative disease
    \item Permanent vision loss and blindness
    \item Reduced quality of life, independence, and driving ability
\end{itemize}

\section*{Prognosis and Outcomes}
With good glucose and blood pressure control and regular screening, most people with diabetes never develop sight-threatening retinopathy. When retinopathy is detected early, laser therapy and modern injectable treatments are highly effective at preserving vision, and blindness from diabetic retinopathy has become largely preventable in Australia. Treatment is more effective the earlier it starts, which is why annual eye screening is so important.

\section*{Prevention and Self-Care}
\begin{itemize}
    \item Keep blood glucose in the target range from diagnosis
    \item Control blood pressure and cholesterol
    \item Attend eye screening at least every one to two years
    \item Report any change in vision promptly
    \item Do not smoke, and limit alcohol
    \item Maintain a healthy weight and stay active
    \item If pregnant, arrange eye assessment, as pregnancy can worsen retinopathy
\end{itemize}

\section*{Sources and Further Reading}
The following reputable sources provide further information for healthcare students and patients seeking reliable, current advice about diabetic retinopathy.

\begin{itemize}
    \item Diabetes Australia. \textit{Diabetic retinopathy: information and screening.} https://www.diabetesaustralia.com.au/
    \item MDFA (Macular Disease Foundation Australia). \textit{Diabetic retinopathy.} https://www.mdfoundation.com.au/
    \item NHS. \textit{Diabetic retinopathy: symptoms and treatment.} https://www.nhs.uk/conditions/diabetic-retinopathy/
    \item Mayo Clinic. \textit{Diabetic retinopathy: diagnosis and treatment.} https://www.mayoclinic.org/
    \item MedlinePlus. \textit{Diabetic retinopathy.} https://medlineplus.gov/
    \item Optometry Australia. \textit{Eye examinations and diabetes.} https://www.optometry.org.au/
\end{itemize}
